\title{Introduction to hyperK\"ahler Geometry}
\author{\textit{Based on a lecture series by }Luca Giovenzana.}
\date{December 2023}
\maketitle

\newcommand{\Def}{\mathop{\mathcal{D\kern -0.5pt e\kern -0.5pt f\kern -1pt}}}
\newcommand{\Moduli}{\ensuremath{\mathscr{M}}}

HyperK\"ahler and K3 manifolds, being interesting objects to study in their own
right, also provide an excellent playground to study ideas in Hodge theory and
deformation/moduli theory. In these lectures we explore some of the
classical theory surrounding the differential and algebraic geometry of these
objects.

\subsection{About} This course on hyperK\"ahler geometry was a part of a Winter School for early career researchers, organised by the UK Algebraic
Geometry Network. The lectures were delivered in-person in the University of Warwick, and were transcribed by Parth Shimpi. These notes have undergone several amendments and are not a verbatim recall of the lectures, therefore discretion is advised when using this material. The notes are available online at
\url{https://pas201.user.srcf.net/documents/2023-ukag-hyperkahler.pdf}. 

All errors and corrections should be communicated to by email to
\url{pas201@srcf.net}.

\section{Basics on K3 surfaces and hyperK\"ahler manifolds}

\begin{definition}
    A compact complex surface \(X\) is a K3 if \(\HH^1(X,\mathscr{O}_X)=0\) and
    \(\Omega^2_X\cong \mathscr{O}_X\).
\end{definition}

\begin{example}
    There are many examples of K3 surfaces, with explicit equations. For
    example, smooth quartic hypersurface \(X\subset \mathbb{P}^3\) is K3.
    Indeed, we have an exact sequence
    \[
        0\to \mathscr{O}_{\mathbb{P}^3}(-3)\to \mathscr{O}_{\mathbb{P}^3}\to
        \mathscr{O}_X\to 0
    \]
    so taking the long exact sequence of cohomology, we find
    \(\HH^1(X,\mathscr{O}_X)=0\). Secondly, by adjunction we see that
    \(K_X=(K_{\mathbb{P}^3}+X)|_X=(-4H+4H)|_X = 0\) so the canonical bundle of
    \(X\) is trivial. A similar argument shows that a smooth surface defined by
    a quadric and a cubic equation in \(\mathbb{P}^4\) is K3. Likewise for three
    quadrics in \(\mathbb{P}^5\).
\end{example}

Generalising the example of a quartic, if \(Y\) is Fano, then a smooth surface
\(X=D_1\cap ...  \cap D_n\) defined as an intersection of divisors is K3 if
\(\sum D_i+K_Y=0\). Thus K3 surfaces often arise from fundamental linear systems
of Fano 3-folds. As another example, if \(\mathbb{V}(f_6)\subset \mathbb{P}^2\)
is a smooth sextic, then \(X=\mathbb{V}(y^2-f_6)\subset \mathbb{P}(1,1,1,3)\) is
a K3 surface in the anticanonical linear system of the weighted projective
space.

\begin{example}
    Another way to obtain K3 surfaces is as quotients of Abelian
    varieties. Consider a complex \(2\)-dimensional torus \(A\), and the map
    \((-1):A\to A\) which sends an element to its inverse under the group law.
    The quotient \(A/\{\pm 1\}\) is singular at precisely the image of the
    \(2\)-torsion subgroup \(A[2]\subset A\) which consists of sixteen points.
    Blowing up at these, we get a K3 surface called the \textit{Kummer
    surface}.
\end{example}

The differential geometry of K3 surfaces is well understood-- in particular they
are all diffeomorphic as smooth \(\mathbb{R}\)-manifolds. Thus there is a unique
lattice (i.e.\ a free abelian group equipped with a symmetric bilinear form)
corresponding to the second integral cohomology of a K3 surface equipped with the
intersection pairing, called the \textit{K3 lattice}. Before stating its
form, we recall some notation from lattice theory-- we write \(U\cong (\mathbb{Z}^2, \left(\begin{smallmatrix} 0&1\\1&0
\end{smallmatrix}\right))\) for the \textit{hyperbolic lattice}. The \(E_8\)
lattice is characterised by its property of being the unique eight-dimensional
lattice that is unimodular (i.e.\ has intersection matrix of determinant \(\pm
1\).) Another lattice that will appear later is \(I_1\), the positive-definite
one-dimensional unimodular lattice. The direct sum of two lattices is always
orthogonal, and if \(\Lambda=(\mathbb{Z}^{\oplus n}, M)\) is a lattice and \(m\)
is an integer, then we define the \textit{twisted lattice} \(\Lambda(m)\) as
\(\mathbb{Z}^{\oplus n}\) equipped with the pairing \(mM\).

\begin{proposition}
    If \(X\) is a K3 surface, then the singular cohomology
    \(\HH^2(X,\mathbb{Z})\) equipped with the intersection pairing is isomorphic
    to the \(22\)-dimensional lattice \(\Lambda_{\textnormal{K3}}=U^{\oplus
    3}\oplus E_8^{\oplus 2}(-1)\). This has signature \((3,19)\).
\end{proposition}

The Torelli theorem for K3 surfaces asserts that the complex structure is
determined by how this sits inside the complex cohomology. This will help
study deformations and moduli of K3 surfaces later.

\subsection{HyperK\"ahler manifolds} K3 surfaces form the two-dimensional
counterpart of a larger family of manifolds with trivial canonical bundles and
symplectic structures, called hyperK\"ahler manifolds.

\begin{definition}
    A complex manifold \((M,I)\) is \textit{K\"ahler} if there exists a metric
    \(g\) compatible with the complex structure \(I\) such that the
    \((1,1)\)-form \(\omega=g(I\cdot -,-)\) is closed. A compact K\"ahler
    manifold \(X\) is \textit{hyperK\"ahler} if it is simply connected (i.e.\
    \(\pi_1(X)=0\)) and the group \(\HH^0(X,\Omega^2_X)\) is generated by a
    holomorphic symplectic form.
\end{definition}

\begin{example}
    \(\mathbb{P}^n_\mathbb{C}\) is a K\"ahler manifold, with K\"ahler form given
    by the Fubini-Study metric \(\omega=\frac{i}{2}\partial \bar\partial
    \log|z|^2\). Thus every projective manifold is also K\"ahler, and not much
    is lost by thinking of the K\"ahler condition as a mild generalisation of
    the projective condition.
\end{example}

Observe that the existence of a symplectic form \(\sigma:T_X\otimes T_X\to
\mathscr{O}_X\) on a hyperK\"ahler manifold \(X\) implies the dimension of \(X\)
is even. Moreover, can use \(\sigma\) to identify \(T_X\) with \(\Hom(T_X,\mathscr{O}_X)=\Omega^1_X\). Thus \(\sigma^{\wedge n}:\mathscr{O}_X\to \Omega^n_X\) gives a trivialisation of the canonical bundle.

Indeed, two dimensional hyperK\"ahler manifolds are precisely K3 surfaces-- by a result of Siu, all K3 surfaces (even the
non-algebraic ones) are K\"ahler. Since all K3 surfaces are diffeomorphic, it
suffices to show any one of them is simply connected and one does this for a
smooth quadric in \(\mathbb{P}^3\) using the hard Lefshetz theorem. Then the
K3 condition implies any holomorphic global function gives a holomorphic
symplectic form generating \(\HH^0(X,\Omega_X^2)\).

In higher dimensions, there aren't many known examples for hyperK\"ahler
manifolds and the ones we know arise from K3s.

\begin{example}[Hilbert schemes of K3 surfaces]
    Let \(S\) be a K3 surface and consider the product \(S^n=S\times...\times
    S\). The symmetric group \(\Sigma_n\) acts on this, and the quotient
    \(S^{(n)}=S^n/\Sigma_n\) is a variety singular along the image of the
    \textit{big diagonal} \(\Delta=\{(x_1,...,x_n)\in S^n\;|\; x_i=x_j \text{
    for some }i\neq j\}\). Beauville proved that there exists a crepant
    resolution \(S^{[n]}\to S^{(n)}\) (called the \textit{Hilbert--Chow
    morphism}), and the variety \(S^{[n]}\) (called the \textit{Hilbert scheme},
    also written \(\text{Hilb}^n(S)\)) parametrises zero-dimensional subschemes
    of length \(n\) in \(S\).

    This is concretely realised in the special case \(n=2\), where we see that
    the Hilbert--Chow morphism \(S^{[2]}\to S^2/\Sigma_2\) sends a length \(2\)
    subscheme to its support (\(Z\mapsto \sum_{P\in
    \text{Supp}Z}\text{length}\mathscr{O}_{Z,P}\cdot P\)) and is the blow-up of the
    singular locus \(\Delta\subset S^2/\Sigma_2\). The exceptional locus
    \(E\cong \mathbb{P}_S(TS)\) is precisely the preimage of points of form
    \(2P\).

    More importantly for us, \(S^{[n]}\) is hyperK\"ahler.
    Indeed it is easy to show the existence of a holomorphic symplectic form--
    since \(S\) is K3, there is a non-trivial element \(\sigma\in
    \HH^0(S,\Omega^2_S)\) and we will use this to construct one on the Hilbert
    scheme. If \(\pi_i:S^n\to S\) denotes the \(i\)th
    projection, then \(\pi_i^\ast \sigma\) is symplectic on \(S^n\) and
    \(\sum_i \pi_i^\ast \sigma\) is invariant under the \(\Sigma_n\)
    action. Thus it descends to a symplectic form on \(S^{(n)}\), and can
    be pulled back via the crepant resolution.
\end{example}

\begin{example}[Generalised Kummer manifolds]
    If \(A\) is a \(2\)-dimensional complex torus, there is a natural map
    \(\Sigma:A^{(n)}\to A\) which applies the group law to a tuple
    \((a_1,...,a_n)\in A^n\). Then considering the composite map
    \(f:A^{[n]}\to A^{(n)}\to A\), the preimage \(\text{Kum}_{n-1}(A)\coloneq
    f^{-1}(0)\) is a \(2(n-1)\) dimensional K\"ahler manifold called the
    \textit{generalised Kummer manifold}.
\end{example}

More examples were provided by O'Grady, who constructed a 6-fold and a 10-fold
as the resolutions of moduli spaces of sheaves on K3 surfaces and abelian
surfaces.

\subsection{Other properties and names} HyperK\"ahler surfaces are so called
because they are very K\"ahler-- we can use the K\"ahler structure \((M,I,g)\)
to find two more complex structures \(J,K\) such that \(g\) is a K\"ahler metric
for all complex structures \(aI+bJ+cK\) parametrised over the unit sphere
\((a,b,c)\in S^2\subset \mathbb{R}^3\).

HyperK\"ahler manifolds are also called \textit{irreducible holomorphically
symplectic manifolds} (IHSMs), because of the result below.

\begin{theorem}[(Beauville--Bogomolov)]
    Let \(X\) be a compact K\"ahler manifold such that \(c_1(X)=0\). Then there
    exists a finite \'etale morphism \(\tilde X\to X\) such that \(\tilde X =
    T\times \prod_j Y_j \times \prod_j Z_j\),  where the factor \(T\) is a
    complex torus, the \(Y_j\) are strict Calabi--Yau (i.e.\ \(K_{Y_j}=0\) and
    \(\HH^0(Y_j, \Omega_{Y_j}^p)=0\) for all \(0<p<\dim Y_j\)) and the \(Z_j\)
    are hyperK\"ahler.
\end{theorem}

\section{Deformations of hyperK\"ahler manifolds}
Because the differential geometry of hyperK\"ahler manifolds is so rigid,
the only thing that can change in a family is the complex structure, which
is largely controlled by the Hodge theory of these manifolds. Being an algebraic
construction, this makes hyperK\"ahler manifolds particularly suited for
studying deformation and moduli problems.

\subsection{Recap: Hodge theory} Recall that a (pure) Hodge structure of weight \(k\) on a finitely generated Abelian group \(H\) is a decomposition \(H\otimes \mathbb{C}=\bigoplus_{p+q=k} H^{p,q}\) such that \(\overline{H^{p,q}}= H^{p,q}\) where \(H\otimes \mathbb{C}\) has a canonical conjugation map given by \(h\otimes \lambda \mapsto h\otimes \bar\lambda\).

If \(X\) is a complex manifold then we have \(\HH^k(X,\mathbb{C})=\HH^k(X,\mathbb{Z})\otimes \mathbb{C}\), and the Hodge decomposition theorem states that \(\HH^k(X,\mathbb{Z})\) has a pure Hodge structure of weight \(k\) whenever \(X\) is a compact K\"ahler manifold. The decomposition factors are given by \(\HH^{p,q}(X)=\HH^q(\Omega_X^p)\), and the Hodge decomposition respects the ring structure on cohomology in that \(\HH^{p,q}(X)\smile \HH^{r,s}(X)\subset \HH^{p+r,q+s}(X)\).

Note also that the subspace \(\HH^{p,q}(X)\oplus \HH^{q,p}(X)\) is stable under
complex conjugation, so considering \(\HH^n(X,\mathbb{R})\) as a subspace of
\(\HH^n(X,\mathbb{C})\) we see that there is an induced decomposition
\(\HH^n(X,\mathbb{R})=\bigoplus_{\begin{smallmatrix}{p+q=n}\\{p<q}\end{smallmatrix}}\HH^{p,q}(X,\mathbb{R})\)
where we write \(\HH^{p,q}(X,\mathbb{R})=\HH^n(X,\mathbb{R})\cap
(\HH^{p,q}(X)\oplus \HH^{q,p}(X))\).

\subsection{The cohomology of hyperK\"ahler manifolds} We saw that the singular
cohomology of a K3 surface is given by the K3 lattice, we will derive a similar
structure if \(X\) is hyperK\"ahler of dimension \(2n\). In this situation, we
have a symplectic form \(\sigma\) generating \(\HH^0(X,\Omega^2_X)\),
normalised so that \(\int_X (\sigma\bar\sigma)^n=1\). This allows us to define a
quadratic form on \(\HH^2(X,\mathbb{Z})\) as
\[
    q_X(\alpha)=\frac{n}{2}
    \int_X \alpha^2(\sigma\bar\sigma)^{n-1}+(1-n)
    \left(\int_X\alpha\sigma^n\bar\sigma^{n-1}\right)
    \left(\int_X\alpha\sigma^{n-1}\bar\sigma^n\right).
\]

Note if \(n=1\) (i.e.\ \(X\) is a K3 surface) then \(q_X=\frac{1}{2}\int_X
\alpha^2\) is proportional to the intersection form. This is true more
generally, in the sense below.

\begin{theorem}
    If \(X\) is a hyperK\"ahler manifold of dimension \(2n\), then there exists
    a \(C\in \mathbb{R}\) such that \(\forall \alpha\in \HH^2(X,\mathbb{Z})\) we
    have \(q_X(\alpha)^n=C\int_X \alpha^{2n}\). In particular, \(q_X\)
    can be rescaled to obtain an integral and primitive non-degenerate form
    \(\HH^2(X,\mathbb{Z})\to \mathbb{Z}\) called the \emph{Beauville--Fujiki
    form} which has signature \((3,\textnormal{b}_2(X)-3)\).
\end{theorem}

We will sketch explicit subspaces on which the form is
positive and negative definite. Since \(\HH^{2,0}(X)\) is spanned by a holomorphic symplectic form \(\sigma\), any \(\alpha\in \HH^2(X,\mathbb{C})\) is of the form \(\alpha=\lambda \sigma + \beta + \mu\bar\sigma\) for \(\beta\in
\HH^{1,1}(X)\) and we then have \(q_X(\alpha)=\lambda\mu + \frac{n}{2}\int_X
\beta^2(\sigma\bar\sigma)^{n-1}.\) In particular \(\HH^{1,1}(X)\) is orthogonal
to \(\HH^{2,0}(X)\oplus \HH^{0,2}(X)\) and \(q_X\) is positive definite on all
forms \(\alpha=\lambda\sigma+\overline{\lambda \sigma}\). The formula also shows
that \(q_X\) is positive definite on any K\"ahler class \([\omega]\in
\HH^{1,1}(X,\mathbb{R})\), so \(q_X\) is positive definite on the three
dimensional subspace \(\HH^{0,2}(X,\mathbb{R})\oplus \mathbb{R}[\omega]\subset
\HH^2(X,\mathbb{R})\).

To understand the orthogonal subspace on which the form is negative
definite, recall that (the cohomology class of) a \(k\)-form \(\alpha\) is
\textit{primitive} if \(\omega^{n-k+1}\cdot \alpha=0\). It follows that a
\(2\)-form \(\alpha\) is primive if and only if \(q_X(\alpha,\omega)=0\), and
hence \(q_X\) is negative definite precisely on the space
\(\HH^{1,1}(X,\mathbb{R})_\text{prim}\) of primitive real (1,1)-forms.

\begin{example}
    If \(X=S^{[n]}\) for \(S\) a K3 surface, then \(\HH^2(X,\mathbb{Z})\cong
    U^{\oplus 3} \oplus E_8(-1)^{\oplus 2} \oplus I_1(2-2n)\) where the lattice
    \(I_1(2-2n)\) is generated by \(\frac{1}{2}[E]\) for \(E\) the exceptional
    divisor of the Hilbert--Chow morphism \(S^{[n]}\to S^{(n)}\).
    On the other hand if \(X=\text{Kum}_n(A)\) for a 2-dimensional complex torus
    \(A\), we have \(\HH^2(X,\mathbb{Z})=U^{\oplus 3}\oplus I_1(-2-2n)\).
\end{example}

\begin{remark}
    For every known example, the form \(q_X:\HH^2(X,\mathbb{Z})\to \mathbb{Z}\)
    is even (i.e.\ has image in \(2\mathbb{Z}\)) but the proof for general
    hyperK\"ahler manifolds remains open.
\end{remark}

\subsection{Deformation families} We can now discuss the main content of this
section, the study of hyperK\"ahler manifolds in families.

\begin{definition}
    A \textit{deformation} of a manifold \(X\) is a smooth proper morphism of
    complex spaces \(\mathcal{X}\to S\), where \(s\in S\) is a distinguished
    point and \(\mathcal{X}|_s\cong X\). The deformation is said to be
    \textit{infinitesimal} if \(S=\Spec \mathbb{C}[\epsilon]/(\epsilon^2)\).

    A deformation \(\mathcal{X}\to S\) is said to be \textit{universal} if for
    any other deformation \(\mathcal{X}'\to S'\ni s'\), there is a unique map
    \((S',s')\to (S,s)\) such that \(\mathcal{X}'\cong \mathcal{X}\times_S
    S'\).
\end{definition}

\begin{remark}
    Note we are working with smooth (not just flat) deformations, and this
    allows a lot of differential geometric tools to come into play. In
    particular by Ehresmann's theorem, whenever \(\mathcal{X}\to S\) is a
    deformation of \(X\), we have an open cover \(S=\bigcup_\alpha U_\alpha\)
    such that \(\mathcal{X}|_{U_\alpha}\) is diffeomorphic to \(X\times
    U_\alpha\). Thus only the complex structure is deformed.
\end{remark}

\begin{theorem}
    The set of infinitesimal deformations of \(X\) up to equivalence is in
    bijection with the space \(\HH^1(X,T_X)\).

    \begin{proof}[Sketch]
        Let \(\mathcal{X}\to \Spec \mathbb{C}[\epsilon]/(\epsilon^2)\) be an
        infinitesimal deformation. By Ehressmann's theorem the normal bundle of
        the central fiber \(X\) is trivial, hence the normal sequence \[0\to
        T_X\to T_{\mathcal{X}}|_X\to \mathscr{O}_X\to 0\] gives an element in
        \(\Ext^1(\mathscr{O}_X, T_X)=\HH^1(X,T_X)\). The inverse map sends an
        element \(v\in \HH^1(X,T_X)\) represented by a 1-cocycle \(\{v_{ij}\in
        \HH^0(U_i\cap U_j, T_X)\}\) to the family \(\mathcal{X}\to
        \Spec\mathbb{C}[\epsilon]/(\epsilon^2)\) constructed by gluing the
        trivial deformations \(U_i\times \Spec
        \mathbb{C}[\epsilon]/(\epsilon^2)\) along the maps \((\text{id}+\epsilon
        v_{ij}):\mathscr{O}_{U_i\cap U_j}[\epsilon]\to\mathscr{O}_{U_i\cap
        U_j}[\epsilon]\).
    \end{proof}
\end{theorem}

Thus the set of infinitesimal deformations is naturally a vector space. It is in
fact the tangent space to the base of the universal deformation \(\mathcal{X}\to
\Def X\ni 0\), if it exists. Indeed in this case infinitesimal deformations are
given by pointed morphisms \(\Spec \mathbb{C}[\epsilon]/(\epsilon^2)\to \Def
X\), and the space of such maps \(\text{Mor}_\ast(\Spec
\mathbb{C}[\epsilon]/(\epsilon^2), \Def X)\) is the tangent space \(T_0\Def
X\).

We will now see that hyperK\"ahler manifolds always have universal deformations
and locally they are simply the affine space \(\mathbb{C}^n\).

\begin{theorem}
    For any complex manifold \(X\), if \(\HH^0(X,T_X)=0\) then a universal
    deformation \(\mathcal{X}\to \Def X\) exists. Moreover for all \(t\in \Def
    X\), this is a universal deformation of \(X_t=\mathcal{X}|_t\).
\end{theorem}

\begin{corollary}
    HyperK\"ahler manifolds have universal deformations.
    \begin{proof}
        If \(X\) is a hyperK\"ahler manifold, then the symplectic form
        \(\sigma:T_X\times T_X\to \mathscr{O}_X\) induces an isomorphism
        \(T_X\cong \Omega_X\) and hence \(\HH^0(X,T_X)=\HH^0(X,\Omega_X)\).
        By Hodge symmetry this has the same dimension as
        \(\HH^1(X,\mathscr{O}_X)\), but this space vanishes for hyperK\"ahler
        manifolds.
    \end{proof}
\end{corollary}

\begin{theorem}[(Bogomolov--Tian--Todorov)]
    If \(X\) is hyperK\"ahler, then the universal deformation space \(\Def X\)
    is smooth (and hence equal to \(\HH^1(X,T_X)\)).
\end{theorem}

\begin{example}
    If \(S\) is a K3 surface, then note that
    \(\text{h}^1(S,T_S)=\text{h}^{1,1}(S)=20\). Can all deformations arise as
    smooth quartics? Note a quartic equation has
    \(\text{h}^0(\mathbb{P}^3, \mathscr{O}_{\mathbb{P}^3}(4))=35\)
    parameters. On the other hand the automorphism group
    \(\text{PGL}(4,\mathbb{C})\) is 15 dimensional, so the space of
    quartics in \(\mathbb{P}^3\) has dimension \(35-1-15=19\). Thus we conclude
    that \(S\) must admit deformation which does not embed in
    \(\mathbb{P}^3\).
\end{example}

\section{Moduli of hyperK\"ahler manifolds}
We have seen that if \(X\) is a hyperK\"ahler manifold then it has a universal
deformation over \(\Def X = \HH^1(X,T_X)\) centered at the origin. Moreover
\(\Def X\) has an open subset \(U_X\) such that the family restricted to \(U_X\)
is a universal deformation of every fiber \(X|_t\) (\(t\in U_X\)). Thus we can
na\"ively consider the set
\[
    \Moduli_X=\{Y\;|\; Y\text{ is hyperK\"ahler and deformation equivalent to
    }X\}
\]
and use the charts from each \(U_Y\subset \Def Y \cong \mathbb{C}^{b_2-2}\) to
give this a complex structure, since if \(Z\in U_Y\) then there is an
open subset of \(U_Z\) identified with a neighbourhood of \(Z\) in \(U_Y\) by
the universality. The same argument also guarantees compatibility at \(Z\in
U_Y\cap U_{Y'}\), since the corresponding restricted universal families
\(\mathcal{Y}|_{U_Y\cap U_{Y'}}\) and \(\mathcal{Y}'|_{U_Y\cap U_{Y'}}\) are
both also universal deformations of \(Z\). Thus the set of all hyperK\"ahler
manifolds \(\Moduli_X\) has the structure of a complex space, but it turns out
this is not Hausdorff and hence isn't a manifold.

To get a well behaved moduli space, we hence parametrise hyperK\"ahler
manifolds with extra data, namely the choice of an isomorphism
\(\HH^2(X,\mathbb{Z})\cong \Lambda_\text{K3}\) and (later) the choice of an
ample bundle. For the rest of the section we will build these ideas for the
specific case of a K3 surface, but everything generalises to hyperK\"ahler
manifolds but with more bookkeeping.

\subsection{Moduli of marked K3 surfaces} As mentioned before, \(K3\) surfaces
have a Torelli theorem which reduces the problem to algebraic data.

\begin{theorem}[(Torelli theorem for K3 surfaces)]
    Two K3 surfaces \(X,X'\)  are isomorphic if and only if there is a Hodge
    isometry \(\HH^2(X,\mathbb{Z})\cong \HH^2(X',\mathbb{Z})\). Moreover,
    such an isometry \(\gamma\) is of the form \(\gamma=f^\ast\) for \(f:X\to
    X'\) if and only if there is a K\"ahler class \(\omega_X\in
    \HH^2(X,\mathbb{R})\) such that \(\gamma(\omega_x)\in \HH^2(X',\mathbb{R})\)
    is K\"ahler.
\end{theorem}

\begin{definition}
    A \textit{marked K3 surface} is a pair \((X,\varphi)\), where \(X\) is a K3
    surface and \(\varphi:\HH^2(X,\mathbb{Z})\to \Lambda_\text{K3}\) is an
    isometry. An isomorphism \((X,\varphi)\to (X',\varphi')\) is an isomorphism
    \(f:X\to X'\) such that \(\varphi\circ f^\ast = \varphi'\).
\end{definition}

By the Torelli theorem, parametrising marked K3 surfaces is equivalent to
parametrising Hodge structures on \(\Lambda_\text{K3}\). Such a structure is
given by a filtration \(0\subset F^0\subset F^1\subset \Lambda_\text{K3}\otimes
\mathbb{C}\) where
\(\dim F^0=1\), \(\dim F^1=21\). Thus to first approximation the parameter space
is in the variety \(\text{Flag}(1,21,\Lambda_\text{K3})\). But this space has
dimension \(21\times 20=420\), which is much larger than the expected dimension of \(\Def X\)!

Indeed not all filtrations can arise-- the Hodge structure in this situation is
completely determined by \(\HH^{2,0}(X)\), since
\(\HH^{0,2}(X)=\overline{\HH^{2,0}}(X)\) and \(\HH^{1,1}=(\HH^{2,0}(X)\oplus
\HH^{0,2}(X))\orth\). Moreover if \(\sigma\) is a generator of
\(\HH^0(X,\Omega^2_X)\), then it satisfies the \textit{Hodge--Riemann bilinear
relations} \((\sigma,\bar\sigma) > 0\), \((\sigma ,\sigma) = 0\). Thus our
desired parameter space is a subset of the Grassmannian
\(\text{Gr}(1,22)=\mathbb{P}(\Lambda_\text{K3}\otimes \mathbb{C})\) given by
\[\Omega_\text{K3}=\{[x]\in \mathbb{P}(\Lambda_\text{K3}\otimes \mathbb{C})
\;|\; (x,x)=0 \text{ and }(x,\bar x)>0\},\] called the \textit{period domain}.
This is a \(20\)-dimensional complex manifold, given by a Euclidean open set in
the subvariety defined by a quadratic form in \(\mathbb{P}^{21}\).

Write \(\Moduli_\text{K3}\) for the set of marked K3 surfaces up to isomorphism. By construction, sending a marked K3 surface \((X,\varphi)\) to the point
\([\varphi(\HH^{2,0}X)]\in \mathbb{P}(\Lambda_{\text{K3}}\otimes \mathbb{C})\)
gives a (surjective!) map \(\Moduli_\textnormal{K3}\to \Omega_\textnormal{K3}\),
called the \textit{period map}. In fact we can give \(\Moduli_\textnormal{K3}\)
an analytic structure by considering families of marked K3 surfaces.

\begin{definition}
    A \textit{family of marked K3 surfaces} is given by a smooth map
    \(\pi:\mathcal{X}\to S\) and an isometry of sheaves \(\Phi:
    \mathbf{R}^2\pi_\ast \underline{\mathbb{Z}}\isoto
    \underline{\Lambda_\text{K3}}\).

    The \textit{period map associated to a family} \((\mathcal{X}\to
    S,\Phi)\) is given by the map \(\mathcal{P}:S\to
    \Omega_{\Lambda_\text{K3}}\), which sends \(s\in S\) to the ``period of
    \(\mathcal{X}|_s\)'', i.e.\ the subspace
    \([\Phi_s(\HH^{2,0}(\mathcal{X}|s))]\).
\end{definition}

Thus the period map for families gives charts on \(\Moduli_\text{K3}\), and the
global period map \(\Moduli_\text{K3}\to \Omega_\text{K3}\) is holomorphic. It
not injective-- if \((X,\varphi)\) and \((X',\varphi')\)  have the same periods,
then theree is an isomorphism \(f:X\to X'\) by the Torelli theorem but
\(\varphi\) may differ from \(f^\ast \circ \varphi'\) by an automorphism of the
lattice (in particular it is possible that such an automorphism that does not
preserve the K\"ahler cone and hence the isometry isn't realised as a pullback).
Thus it makes sense to quotient by the action of the orthogonal group
\(\text{O}(\Lambda_\text{K3})\) to obtain a map
\[
    \Moduli_\text{K3}/\text{O}(\Lambda_\text{K3})\longrightarrow
    \Omega_\text{K3}/\text{O}(\Lambda_\text{K3}).
\]
By the Torelli theorem this map is an isomorphism of complex spaces, and the
points of this orbit space are isomorphism classes of K3 surfaces.

The connected components of \(\Moduli_\text{K3}/\text{O}(\Lambda_\text{K3})\)
are the moduli spaces of K3 surfaces that belong to a particular deformation
family, this is precisely the construction that came from considering universal
deformations at the start of this section. In this context, the non--Hausdorff
behaviour can be seen as coming from the fact that the action of
\(\text{O}(\Lambda_\text{K3})\) is far from being properly discontinuous.

\subsection{Algebraicity of K3 surfaces} The moduli space constructed above
is ill-behaved partly because we are considering `too many' surfaces-- most K3
surfaces in the moduli space are not algebraic.

\begin{theorem}[(Kodaira)]
    A complex manifold \(X\) is projective if and only if
    \(\HH^2(X,\mathbb{Z})\) contains a K\"ahler class.
\end{theorem}

Indeed by the Lefshetz \((1,1)\)-theorem we know that every integral
\((1,1)\)-class (i.e.\ element of \(\HH^2(X,\mathbb{Z})\cap \HH^{1,1}(X)\)) is
the Chern class of some holomorphic line bundle, and Kodaira's theorem shows
that the classes with positive representatives (i.e.\ K\"ahler classes) come
from ample line bundles.

\begin{corollary}
    A very general\footnote{We use `very general' to mean `outside the union
    of countably many Zariski-closed subvarieties'.} K3 surface is not
    projective.
    \begin{proof}
        If \(X\) is a projective K3 surface, then \(\HH^{1,1}(X)\cap
        \HH^2(X,\mathbb{Z})\) is necessarily non-zero and hence the subspace
        \(\HH^{2,0}(X)\) (which determines the complex structure on \(X\)) is
        orthogonal to some \(\delta \in \HH^2(X,\mathbb{Z})\). It follows that
        in \(\Omega_\text{K3}\), the locus of marked projective K3 surfaces is
        contained in the closed hyperplane arrangement
        \(\Omega_\text{K3}\cap \left(\bigcup_{\delta\in
        \Lambda_\text{K3}}\delta\orth\right)\). The complement \(U\) is dense
        and open, and any K3 surface with period in \(U\) is not projective.
    \end{proof}
\end{corollary}

Thus to get better behaved moduli spaces, we consider polarised K3 surfaces.

\begin{definition}
    A \textit{polarised K3 surface} is a pair \((X,L)\) where \(X\) is a K3
    surface and \(L\) is a primitive (i.e.\ not of the form \(M^{\otimes (>1)}\))
    ample line bundle. The \textit{degree} of the polarised K3 surface \((X,L)\)
    is equal to \(\text{deg}L=L^2\).
\end{definition}

\subsection{Moduli of polarised K3 surfaces} Note the degree of a polarised K3
surface \(X\) is always even and positive, say equal to \(2d\). Now in this
situation \(L^{\otimes 3}\) is very ample, giving an embedding \(X\into
\mathbb{P}^{9d+1}\) such that \(X\) has Hilbert polynomial \(f(t)=9d t^2+2\).
Thus the straightforward thing to do is to realise that polarised K3 surfaces of
degree \(2d\) lie in an open subset of the Hilbert scheme
\(\text{Hilb}^{9dt^2+2}(\mathbb{P}^{9d+1})\), and thus the required moduli space
of K3 surfaces is obtained as a GIT quotient of this locus by the
\(\text{PGL}(9d+1,\mathbb{C})\) action. In practice, however, checking for
semistability is hard.

This is circumvented by constructing the moduli space explicitly from that of
marked K3 surfaces and using a stronger form of the Torelli theorem. We will
provide a sketch.

By standard lattice theory (Eichler's criterion), up to the action of
\(\text{O}(\Lambda_\text{K3})\) there is a unique primitive element \(\ell\in
\Lambda_\text{K3}\) satisfying \(\ell^2=2d\) and hence we can choose a marking
\(\varphi:\HH^2(X,\mathbb{Z})\to \Lambda_\text{K3}\) with
\(\varphi(\text{c}_1L)=\ell\). Then its reasonable to consider the set
\[
    \Moduli_{2d}= \{(X,\varphi,L)\;|\; (X,\varphi)\text{ is a
    marked K3 surface, }(X,L)\text{ is polarised, and }\varphi[L]=\ell\}.
\]
Now if \((X,\varphi,L)\in \Moduli_{2d}\), then clearly the period point
\([\varphi(\HH^{2,0}X)]\in \mathbb{P}(\Lambda_\text{K3}\otimes \mathbb{C})\)
lies in the hyperplane \(\{[x]\;|\; x \perp \ell\}\), which we can write as the
projective space \(\mathbb{P}(\Lambda_\text{2d}\otimes \mathbb{C})\) for
\(\Lambda_\text{2d}=\ell\orth \cong U^{\oplus 2}\oplus E_8(-1)^{\oplus 2}\oplus
I_1(-2d)\). Thus the natural period domain to consider is
\[
    \Omega_{2d}=\{[x]\in \mathbb{P}(\Lambda_{2d}\otimes \mathbb{C})\;|\;
    (x,x)=0,\; (x,\bar x)>0\}.
\]
Turns out this has two irreducible components which are interchanged by complex
conjugation, so we take the orientation preserving component \(\Omega_{2d}^+\).
This is an example of a \textit{Seigel upper half plane}, a family of spaces
forming higher dimensional analogue of the upper half plane \(\mathbb{H}\).
There is a natural period map \(\Moduli_{2d}\to \Omega_{2d}\), and considering
families of marked polarised K3 surfaces (defined in the obvious way) give
\(\Moduli_{2d}\) a holomorphic structure.

Again, to forget the marking we consider the action of
\(\text{Stab}(\ell)\subset O(\Lambda_{K3})\) on \(\Moduli_{2d}\) and that of
\(\text{Stab}(\ell)^+\) on \(\Omega_{2d}^+\) (where we take the subgroup
preserving the chosen component).  It is convenient to write
\(\Omega_{2d}/\text{Stab}(\ell)\) for \(\Omega_{2d}^+/\text{Stab}(\ell)^+\),
this is an irreducible normal quasi-projective variety by a theorem of
Bailey--Borel. Then one way to state the \textit{strong Torelli theorem for K3
surfaces} is that the induced period map \(\Moduli_{2d}/\text{Stab}(\ell)\to
\Omega_{2d}/\text{Stab}(\ell)\) is an open immersion. This gives us the moduli
of polarised K3 surfaces as a 19 dimensional quasi-projective variety.

Note the difference from the moduli space of marked K3 surfaces we constructed
previously-- the period map here not surjective. To compute the complement, recall that the Nakai--Kleiman criterion states a
line bundle \(L\) on a surface \(X\) is ample if and only if \(L^2>0\) and
\(L\cdot C>0\) for all curves \( C\subset X\). When \(X\) is K3, it suffices
to test the second condition only on curves with self-intersection \(-2\), i.e.\
the \textit{roots} of \(\Lambda_{2d}\). Writing \(\Delta\subset \Lambda_{2d}\)
for the set of roots, we thus see that the complement
\((\Omega_{2d}/\text{Stab}(\ell)) \setminus (\Moduli_{2d}/\text{Stab}(\ell))\)
is given by the root hyperplanes
\(\bigcup_{\delta \in \Delta}\{[x]\;|\; x\perp \delta\}\).

\begin{remark}
    For general hyperK\"ahlers there is a very similar picture that realises all
    deformation families. The slight subtlety is that the polarisation \(\ell\in
    \Lambda_X\) is no longer uniquely determined, and hence it is also not
    obvious what the replacement for the group action should be.
    Turns out the solution is to consider the group of Hodge isometries arising
    as monodromies, but that is beyond the scope of these notes.
\end{remark}

\section{Cubic 4-folds}

In this section we will use some of the techniques discussed to address
rationality problems for cubic 4-folds. To begin, note that the locus of smooth cubic 4-folds forms an open set \(V\) in the linear system \(|\mathscr{O}_{\mathbb{P}^5}(3)|\cong \mathbb{P}^{55}\). Since the automorphism group \(\text{PGL}(6,\mathbb{C})\) is \(35\)-dimensional, we see that the quotient \(\mathcal{C}=V\sslash \text{PGL}(6,\mathbb{C})\) is a \(20\)-dimensional quasi-projective variety, the \textit{moduli space of cubic
4-folds}. Conjecturally rational cubics form a codimension one locus related
intricately with the Hodge theoretic and derived=categorical behaviour of K3
surfaces.

Before proceeding, we first recall some facts about the cohomology of cubic 4-folds. If \(Y\subset \mathbb{P}^5\) is a cubic hypersuface, its Hodge diamond looks like
\[\begin{tikzcd}[column sep=0em, row sep = 0em]
    &&&& 1 &&&& \\
    &&& 0 && 0 &&&\\
    && 0 && 1 && 0 && \\
    &0 &&0 &&0 &&0&\\
    0&&1&&21&&1&&0.
\end{tikzcd} \]
The most interesting bit is the middle cohomology lattice
\(\HH^4(Y,\mathbb{Z})\cong I_1^{\oplus 21}\oplus I_1(-1)^{\oplus 2}\). We can pick a generator \(h\in \HH^2(X,\mathbb{Z})\) associated
(via the Chern class map) to a hyperplane section, then the class \(h^2\in
\HH^4(X,\mathbb{Z})\) has self intersection \(3\). This allows us to define
the \textit{primitive cohomology lattice}
\(\HH^4(Y,\mathbb{Z})_\text{prim}=(\hh^2)\orth \cong E_8^{\oplus 2}\oplus U^{\oplus 2}\oplus
\left(\begin{smallmatrix}2&1\\1&2\end{smallmatrix}\right)\).

Note the similarity with the K3 lattice-- in fact the primitive cohomology of
certain cubic 4-folds can be identified with a subspace of the cohomology of an
associated K3 surface, so that these cubics `look like K3 surfaces' to Hodge
theory. Conjecturally these are precisely the rational cubics. We will sketch this with examples.

\subsection{Explicit constructions} We look at two examples of
rational cubic 4-folds and how K3 surfaces naturally arise in this setting.

\begin{example}[Cubics containing two planes]
    Consider the cubic polynomials \(g=u^2x+v^2y+w^2z\) and
    \(h=ux^2+vy^2+wz^2\), and the associated hypersurface
    \(Y=\mathbb{V}(g-h)\subset \mathbb{P}^5\). This is smooth, and one can
    consider the contained planes \(\Pi_1=\mathbb{V}(u,w,v)\) and
    \(\Pi_2=\mathbb{V}(x,y,z)\). Given
    two general points \(p_1\in \Pi_1\), \(p_2\in \Pi_2\), the line through them
    \[\ell_{p_1,p_2}=\{s\cdot p_1+ t\cdot p_2\;|\; [s:t]\in \mathbb{P}^1\}\]
    intersects \(Y\) in three points \(\{p_1,p_2,q\}\) by Bezout's
    theorem; this defines a rational map \(\Pi_1\times \Pi_2\dashrightarrow Y\)
    given by \((p_1,p_2)\mapsto q\). It is a birational map-- given any point
    \(q\in \mathbb{P}^5\setminus (\Pi_1\cup \Pi_2)\), there is a unique line
    passing through \(q\) that meets \(\Pi_1\) and \(\Pi_2\). Thus \(Y\) is
    birational to \(\mathbb{P}^2\times \mathbb{P}^2\), i.e.\ it is rational.

    The given map \(\Pi_1\times \Pi_2\dashrightarrow Y\)  is undefined whenever
    \(\ell_{p_1,p_2}\) lies  inside \(Y\), i.e.\ on \(S=\mathbb{V}(g,h)\cap
    \Pi_1\cap \Pi_2\). Since \(S\) is a
    \((1,2),(2,1)\) complete intersection in \(\mathbb{P}^2\times
    \mathbb{P}^2\), it is a K3 surface.
\end{example}

\begin{example}[Pfaffian cubics]
    Fix a vector space \(V\cong \mathbb{C}^6\) and consider the Grassmannian
    \(G=\text{Gr}(2,V)\) parametrising lines in \(\mathbb{P}(V)\). This is
    8-dimensional, sits as a degree 14 subvariety inside
    \(\mathbb{P}(\Exter^2V)\) via the Pl\"ucker embedding, and is Fano with
    canonical bundle \(\mathscr{O}_G(-6)\). Thus for a general 9-dimensional
    linear subspace \(L\subset \Exter^2V\), the variety \(S=G\cap
    \mathbb{P}(L)\) is the complete intersection of \(G\) with six hyperplanes
    (corresponding to a basis of \(L\orth\subset \Exter^2V^\ast\)), and
    adjunction formula shows \(S\) is a K3 surface. In fact all general
    polarised K3 surfaces \((S,H)\) of degree \(H^2=14\) arise in this fashion
    \cite[see][\S 3]{mukaiCurvesK3Surfaces}.

    There is a natural way to produce cubic 4-folds in this setting. Note that
    the projective space \(\mathbb{P}(\Exter^2V^\ast)\) parametrises alternating
    forms \(V\times V\to \mathbb{C}\) up to scaling. Now the
    determinant of a \(6\times6\) skew-symmetric matrix \(M\) has form
    \(\det M = \text{Pf}(M)^2\) for a homogeneous cubic polynomial
    \(\text{Pf}\in \mathbb{C}[M_{ij}\;|\; 1\leq i<j\leq 6]\), so the vanishing
    locus \(\mathbb{V}(\text{Pf})\subset \mathbb{P}(\Exter^2V^\ast)\)
    parametrises degenerate forms (i.e.\ forms of rank \(\leq 4\)). On the other
    hand the Grassmannian \(G^\ast = \text{Gr}(2,V^\ast)\) parametrises forms
    of rank \(2\), thus giving us a natural stratification
    \(G^\ast\subset  \mathbb{V}(\text{Pf}) \subset \mathbb{P}(\Exter^2V^\ast).
    \) If the choice of \(L\) was sufficiently general, then \(L\orth\subset
    \Exter^2V^\ast\) is disjoint from \(G^\ast\) and is not
    contained in \(\mathbb{V}(\text{Pf})\). Thus \(\mathbb{P}(L\orth)\) is a
    \(\mathbb{P}^5\) whose points correspond to skew-forms \(V\times V\to
    \mathbb{C}\) of rank \(\geq 4\). The locus of degenerate forms in
    \(\mathbb{P}(L\orth)\) is then
    given by a cubic hypersurface \(Y=\mathbb{P}(L\orth)\cap
    \mathbb{V}(\text{Pf})\), called the \textit{Pfaffian cubic} associated to
    \(L\). If explicit coordinates are fixed, then the choice of \(L\orth\)
    corresponds to the choice of a \(6\times 6\) skew-symmetric matrix with
    entries given by linear forms in \(\mathbb{C}[u,v,w,x,y,z]\), and \(Y\) cuts
    out the locus where this matrix is degenerate.
\end{example}

\begin{exercise}
    Show that \(Y\) is rational, by considering a general 5-dimensional subspace
    \(W\subset V\) and mapping \([\varphi]\in Y\) to \([W\cap \ker\varphi]\in
    \mathbb{P}(W)\). The birational map \(\mathbb{P}^4\dashrightarrow S\) is in
    fact a blow-up in a surface birational to \(S\) \cite[see][\S
    4.1.3]{hassettSpecialCubicFourfolds}.
\end{exercise}

Thus to the data of \(L\subset \Exter^2V\), we have associated a cubic 4-fold
\(Y\) parametrising non-degenerate alternating forms in \(L\orth\) and a K3
surface \(S\) parametrising the 2-planes \(P\subset V\)  such that
\(\varphi|_P=0\) for every \(\varphi\in L\orth\). There is a surprising
connection between the geometry of the two, first explored by Beauville and
Donagi in \cite{beauvilleVarieteDroitesHypersurface}. They were studying the
so-called \textit{Fano variety of lines} \(F(Y)\), which parametrises
lines contained in \(Y\). Since \(Y\) lies in \(\mathbb{P}(L\orth)\), the
variety \(F(Y)\) naturally sits in the Grassmannian
\(\text{Gr}(2,L\orth)\).

\begin{theorem}[(Beauville--Donagi)]
    The Fano variety of lines \(F(Y)\) is isomorphic to the Hilbert scheme
    \(S^{[2]}\).

    \begin{proof}[Sketch]
        Given two distinct points in \(S\), we can realise them as \(2\)-planes
        \(P,Q\subset V\). By the generality hypothesis \(P+Q\)
        is a 4-plane, and we consider the linear subspace \(\ell=\{[\varphi]\in
        \mathbb{P}(L\orth)\;:\; \varphi|_{P+Q}=0\}\). This is 1-dimensional, and
        if \([\varphi]\in \ell\) then \(\varphi|_P=\varphi|_Q=0\) so
        \(\varphi\) is necessarily degenerate. Thus we have a line \(\ell\subset
        Y\), and our map \(S^{[2]}\to F(Y)\) maps \(\{P,Q\}\mapsto \ell\).
    \end{proof}
\end{theorem}

\begin{remark}
    Despite its name, \(F(Y)\) is not a Fano variety. Indeed the theorem shows
    it is hyperK\"ahler, and the name is simply a reflection of the fact that
    this variety was first extensively studied by Fano.
\end{remark}

\subsection{Hodge theory of cubic 4-folds} If \(Y\subset \mathbb{P}^5\) is a
cubic 4-fold, we can analogously construct the Fano variety of lines
\(F(Y)\subset \text{Gr}(2,6)\), this is a smooth hyperK\"ahler 4-fold. There is
a universal line \(\mathscr{L}\subset F(Y)\times Y\) parametrising pairs
\((\ell,y)\) such that \(y\in \ell\). The natural map \(q:\mathscr{L}\to F(Y)\)
is a \(\mathbb{P}^1\)--bundle, and the fibers of the map \(p:\mathscr{L}\to Y\)
are lines in \(\text{Gr}(2,6)\). There is an induced morphism on cohomology
\[
    \alpha=q_\ast\circ p^\ast : \HH^4(Y,\mathbb{Z})\to \HH^2(F_1(Y), \mathbb{Z})
\]
called the \textit{Abel--Jacobi map}, which sends the class \(h^2\) to
\(g=[\mathscr{O}_{F(Y)}(1)]\) where the hyperplane class on \(F(Y)\) is induced
from the Pl\"ucker embedding of the ambient Grassmannian. This map in fact
preserves the Hodge structures, up to a twist.

\begin{theorem}
    The map
    \(\alpha:\HH^4(Y,\mathbb{Z})\to \HH^2(F(Y), \mathbb{Z})(-1)\) is an
    isomorphism of Hodge structures, and induces a Hodge isometry \((h^2)\orth
    \cong g\orth\).
\end{theorem}

Now \cite{charlesRemarkTorelliTheorem} shows that \(F(Y)\) in fact
determines \(Y\). Since hyperK\"ahler manifolds have a Torelli theorem, we can
use this construction to obtain an analogous one for cubic 4-folds.

\begin{corollary}
    Two cubic 4-folds are isomorphic if and only if their primitive
    cohomology lattices are isometric as Hodge structures.
\end{corollary}

Thus fixing the lattice \(M=U^{\oplus 2}\oplus E_8^{\oplus 2}\oplus
\left(\begin{smallmatrix}2&1\\1&2\end{smallmatrix}\right)\), we can consider the
manifold \(\Omega_M\) parametrising compatible Hodge structures on \(M\). As
before, there is a period map \(\mathcal{P}:\mathcal{C}\to \Omega_M/\text{O}(M)\) from the moduli space
\(\mathcal{C}\) of cubic 4-folds. The Torelli theorem says this period map is an
open immersion.

\subsection{Rationality problems and special cubics} In particular, the period
map can be used to show that a very general cubic 4-fold has
\(\HH^{2,2}(X,\mathbb{Z})\cong \mathbb{Z}\cdot h^2\).

\begin{definition}
    Say a cubic 4-fold \(Y\) is \textit{special} if there is a surface
    \(S\subset Y\) with \([S]\notin \mathbb{Z}\cdot h^2\subseteq
    \HH^{2,2}(X,\mathbb{Z})\).

    For such \(Y\), we say a \textit{labelling} is a saturated rank \(2\)
    sublattice \(K\subset \HH^{2,2}(X,\mathbb{Z})\) containing \(\mathbb{Z}\cdot
    h^2\). The \textit{discriminant} of the labelling is the determinant of the
    Gram matrix of the lattice.

    Say a polarised K3 surface \((X,L)\) is \textit{associated to} a labelled
    cubic 4-fold \((Y,K)\) if \(K\orth \subset \HH^4(Y,\mathbb{Z})\) is
    isometric as a Hodge structure to \([L]\orth\subset \HH^2(X,\mathbb{Z})\).
\end{definition}

These were extensively studied in \cite{hassettSpecialCubicFourfolds}, and
we summarise some of the results.

\begin{theorem}[(Hasset)]
    For \(d\geq 0\), consider the space \(\mathcal{C}_d\subset \mathcal{C}\) of
    smooth cubic 4-folds that admit a labelling of discriminant \(d\). Then
    \(\mathcal{C}_d\) is an irreducible divisor if and only if \(d>6\) and
    \(d\equiv 0,2\pmod 3\)
\end{theorem}

\begin{theorem}[(Hasset)]
    A labelled cubic 4-fold \((Y,K)\) has an associated K3 surface if and
    only if its discriminant \(d\) is positive, \(4,9\nmid d\), and all odd
    prime factors of \(d\) are of the form \(3n+1\).
\end{theorem}

Conjecturally, rationality is equivalent to having an associated K3 surface.

\textbf{Conjecture.} (Hasset) A cubic 4-fold \(Y\) is rational if and only if
\(Y\in \mathcal{C}_d\) for \(d\) satisfying the conditions of the theorem
above. In particular, a very general cubic 4-fold is irrational.

On the other hand, by considering \textit{non-commutative K3 surfaces} arising
as Kuznetsov components in \(\Db Y\), Kuznetsov conjectures another criterion
for rationality.

\textbf{Conjecture.} (Kuznetsov) A cubic 4-fold \(Y\) is rational if and only if
the Kuznetsov component associated to the exceptional collection
\((\mathscr{O}_Y,\mathscr{O}_Y(1),\mathscr{O}_Y(2))\) is equivalent to the
derived category of a K3 surface.
